\newcommand{\mail}{\includegraphics[width=0.2\textwidth]{figures/mail.pdf}}
\tikzstyle{arr}=[line width=1mm,draw=black,-triangle 45]
\tabulinesep=4pt

\headlineframe{Versionskontrolle}
\begin{frame}[c]{Was ist das?}
  \begin{itemize}
    \item Verwaltung von Versionen
    \item Speicherung der \enquote{Geschichte} eines Projekts
    \item Es ist jederzeit möglich auf eine ältere Version zurückzukehren
    \item Es ist möglich, sich die Unterschiede zwischen Versionen anzeigen zu lassen
    \item Backup
  \end{itemize}
  \vspace{1cm}
  \begin{center}
    \Large\alert{Wichtige Voraussetzungen für korrektes wissenschaftliches Arbeiten, \\ auch wenn man alleine arbeitet}
  \end{center}
\end{frame}

\headlineframe{Wie arbeitet man am besten an einem Protokoll zusammen?}

\headlineframe{Idee: Austausch über Mails}

\begin{frame}{Mails: Probleme}
  \centering
  \begin{tikzpicture}
    \node (alice) at (0, 0) {\includegraphics[width=0.2\textwidth]{figures/text.pdf}};
    \node (bob) at (8, 0) {\includegraphics[width=0.2\textwidth]{figures/text.pdf}};
    \draw[arr] (alice) -- node {\mail} (bob);
  \end{tikzpicture}
  \begin{itemize}
    \item Risiko, dass Änderungen vergessen werden, ist groß
    \item Bei jedem Abgleich muss jemand anders aktiv werden
      \begin{itemize}
        \item Stört
        \item Es kommt zu Verzögerungen
      \end{itemize}
  \end{itemize}
  \textbf{\Large Fazit: Eine sehr unbequeme / riskante Lösung}
\end{frame}

\headlineframe{Idee: Austausch über Dropbox}

\begin{frame}{Dropbox: Probleme}
  \centering
  \begin{tikzpicture}
    \node (alice) at (0, 0) {\includegraphics[width=0.2\textwidth]{figures/text.pdf}};
    \node[align=center,fill=blue!20,rounded corners=1em] (dropbox) at (4, 0.5) {
      \includegraphics[width=0.1\textwidth]{logos/dropbox.pdf}\\
      \includegraphics[width=0.2\textwidth]{figures/text.pdf}
    };
    \node (bob) at (8, 0) {\includegraphics[width=0.2\textwidth]{figures/text.pdf}};
    \draw[arr,triangle 45-triangle 45] (0.8, 0) -- (dropbox);
    \draw[arr,triangle 45-triangle 45] (dropbox) -- (7.2, 0);
  \end{tikzpicture}

  \begin{itemize}
    \item Man merkt nichts von Änderungen der Anderen
    \item Gleichzeitige Änderungen führen zu \enquote{In Konflikt stehende Kopie}-Dateien
    \item Änderungen werden nicht zusammengeführt
    \item Keine echte Historie des Projekts
  \end{itemize}
  \textbf{\Large Fazit: Besser, aber hat deutliche Probleme}
\end{frame}

\headlineframe{Lösung: Änderungen verwalten mit \texttt{git}}

\begin{frame}
    \centering
    \includegraphics[width=0.7\textwidth]{logos/git.pdf}

    \vspace{1em}

    \begin{itemize}
      \item Ein Versionskontrollsystem
      \item Ursprünglich entwickelt, um den Programmcode des Linux-Kernels zu verwalten (Linus Torvalds)
      \item Hat sich gegenüber ähnlichen Programmen (SVN, mercurial) durchgesetzt
      \item Wird in der Regel über die Kommandozeile benutzt
    \end{itemize}
\end{frame}

\begin{frame}{Was bringt \texttt{git} für Vorteile?}
  \begin{itemize}
    \item Arbeit wird für andere sichtbar protokolliert
    \item Erlaubt Zurückspringen an einen früheren Zeitpunkt
    \item Kann die meisten Änderungen automatisch zusammenfügen
    \item Wirkt nebenbei auch als Backup
  \end{itemize}
  Einzige Herausforderung: Man muss lernen, damit umzugehen
\end{frame}

\headlineframe{Das Repository}
\begin{frame}{Zentrales Konzept: Das Repository}
  \begin{itemize}
    \item Erzeugen mit \texttt{git init}
    \item Damit wird der aktuelle Ordner zu einem Repository
  \end{itemize}
  \vspace{3em}
  \centering
  \begin{tikzpicture}[
      line width=1.5,
      gitstep/.style={
        draw,
        rounded corners,
        thick,
        minimum width=4cm,
        minimum height=1.2cm,
      },
    ]
    \node (wd) at (0, 0) [gitstep, fill=red!20, visible on=<1->] {Working directory};
    \node (idx) [gitstep, fill=yellow!20, visible on=<2->, below=0.4cm of wd] {Staging};
    \node (hist) [gitstep, fill=green!20, visible on=<3->, below=0.4cm of idx] {History};
    \draw[thick,->, visible on=<4->] (wd.east) to[out=0,in=0] node[right] {\texttt{git add}} (idx.east);
    \draw[thick,->, visible on=<4->] (idx.west) to[out=180,in=180] node[left] {\texttt{git commit}} (hist.west);
  \end{tikzpicture}
\end{frame}
\begin{frame}{Zentrales Konzept: Das Repository}
  \centering
  \begin{tikzpicture}[
      line width=1.5,
      gitstep/.style={
        draw,
        rounded corners,
        thick,
        minimum width=4cm,
        minimum height=1.2cm,
      },
    ]
    \node (wd) at (0, 0) [gitstep, fill=red!20, visible on=<1->] {Working directory};
    \node (idx) [gitstep, fill=yellow!20, below=0.4cm of wd] {Staging};
    \node (hist) [gitstep, fill=green!20, below=0.4cm of idx] {History};
    \node [right=1cm of wd, text width=0.5\textwidth, align=flush left] {Aktuelles Arbeitsverzeichnis, Inhalt des Ordners im Dateisystem.};
    \node [right=1cm of idx, text width=0.5\textwidth, align=flush left] {Änderungen, die für den nächsten commit vorgemerkt sind.};
    \node [right=1cm of hist, text width=0.5\textwidth, align=flush left] {Gespeicherte Historie des Projekts. Alle jemals gemachten Änderungen. Ein Baum von Commits.};
  \end{tikzpicture}
\end{frame}

\begin{frame}{Remotes}
  Remotes sind zentrale Stellen, z.\,B. Server auf denen die History gespeichert wird.
  \begin{center}
  \begin{tikzpicture}[line width=1.5]
    \node (hist) [draw,rounded corners,thick,minimum width=4cm,minimum height=1.2cm,fill=green!20] {History};
    \node (remote) [below=0.4cm of hist,draw,rounded corners,thick,minimum width=4cm,minimum height=1.2cm,fill=blue!20] {Remote};
    \draw[thick,->] (hist.east) to[out=0,in=0] node[right] {\texttt{git push}} (remote.east);
    \draw[thick,->] (remote.west) to[out=180,in=180] node[left] {\texttt{git pull}} (hist.west);
  \end{tikzpicture}
  \end{center}
\end{frame}


\begin{frame}{History}
  \only<1>{
    \begin{tikzpicture}
      \graph [
        nodes={
          blue!70!black,
          node font=\ttfamily,
        },
      ]{
        "Erstabgabe" [visible on=<4->, green!60!black, x=4];
        a <- b <- c <- d <- master [vertexDarkRed];
        f[visible on=<2>] ->[visble on=<2>] -> a;
        "Erstabgabe"[visible on=<2>] ->[visible on=<3>] c;
      };
    \end{tikzpicture}
  }
  \only<2>{
    \begin{tikzpicture}
      \graph [
        nodes={
          blue!70!black,
          node font=\ttfamily,
        },
      ]{
        "Erstabgabe" [visible on=<4->, green!60!black, x=4];
        a <- b <- c <- {
          d <- master [vertexDarkRed],
          e <- foo [vertexDarkRed]
        };
        "Erstabgabe"[visible on=<3>] ->[visible on=<3>] c;
      };
    \end{tikzpicture}
  }
  \only<3->{
    \begin{tikzpicture}
      \graph [
        nodes={
          blue!70!black,
          node font=\ttfamily,
        },
      ] {
        "Erstabgabe" [visible on=<4->, green!60!black, x=4],
        a <- b <- c <- d <- {e <- f, g [x=0.5]} <- h <- {i <-[visible on=<2->] master [vertexDarkRed, visible on=<3->], j <-[visible on=<2->] foo [vertexDarkRed, visible on=<3->]};
        "Erstabgabe" ->[visible on=<4->] e ;
      };
    \end{tikzpicture}
  }

  \vspace{1em}
  \begin{itemize}
    \item<1-> \textcolor{blue}{Commit}: Zustand/Inhalt des Arbeitsverzeichnisses zu einem Zeitpunkt
      \begin{itemize}
        \item Enthält Commit-Message (Beschreibung der Änderungen)
        \item Wird über einen Hash-Code identifiziert
        \item Zeigt immer auf seine(n) Vorgänger
      \end{itemize}
    \item<2-> \textcolor{vertexDarkRed}{Branch}: benannter Zeiger auf einen Commit
      \begin{itemize}
        \item Entwicklungszweig
        \item Im Praktikum reicht bereits die Standard-Branch: \texttt{master}
        \item Wandert weiter
      \end{itemize}
    \item<4-> \textcolor{green!60!black}{Tag}: unveränderbarer Zeiger auf einen Commit
      \begin{itemize}
        \item Wichtiges Ereignis, z.B. veröffentlichte Version
      \end{itemize}
  \end{itemize}
\end{frame}

\begin{frame}{Typischer Arbeitsablauf}
  \begin{enumerate}
    \item Neues Repo? Repository erzeugen oder klonen: \hfill\texttt{git init}, \texttt{git clone} \\
      Repo schon da? Änderungen herunterladen: \hfill\texttt{git pull}
    \item Arbeiten
      \begin{enumerate}
        \item Dateien bearbeiten und testen
        \item Änderungen vorbereiten: \hfill\texttt{git add}
        \item Änderungen als \emph{commit} speichern: \hfill\texttt{git commit}
      \end{enumerate}
    \item Commits anderer herunterladen und integrieren: \hfill\texttt{git pull}
    \item Eigene Commits hochladen: \hfill\texttt{git push}
  \end{enumerate}
\end{frame}

\begin{frame}{\texttt{git init}, \texttt{git clone}}
  \begin{tabu}{>{\ttfamily}l X[,L]}
    git init               & initialisiert ein \texttt{git}-Repo im jetzigen Verzeichnis \\
    git clone \textit{url} & klont das Repo aus \texttt{\textit{url}} \\
    rm -rf .git            & löscht alle Spuren von \texttt{git} aus dem Repo
  \end{tabu}
\end{frame}

\begin{frame}{\texttt{git status}, \texttt{git log}}
  \begin{tabu}{>{\ttfamily}l X[,L]}
    git status & zeigt Status des Repos (welche Dateien sind neu, gelöscht, verschoben, bearbeitet) \\
    git log    & listet Commits in aktuellem Branch
  \end{tabu}
\end{frame}

\begin{frame}{\texttt{git add}, \texttt{git mv}, \texttt{git rm}, \texttt{git reset}}
  \begin{tabu}{>{\ttfamily}l X[,L]}
    git add \textit{file} … & fügt Dateien/Verzeichnisse zum Staging-Bereich hinzu \\
    git add -p …            & fügt Teile einer Datei zum Staging-Bereich hinzu \\
    git mv                  & wie \texttt{mv} (automatisch in Staging)\\
    git rm                  & wie \texttt{rm} (automatisch in Staging) \\
    git reset \textit{file} & entfernt Dateien/Verzeichnisse aus Staging
  \end{tabu}
\end{frame}

\begin{frame}{\texttt{git diff}}
  \begin{tabu}{>{\ttfamily}l X[,L]}
    git diff                                   & zeigt Unterschiede zwischen Staging und Arbeitsverzeichnis \\
    git diff --staged                          & zeigt Unterschiede zwischen letzten Commit und Staging \\
    git diff \textit{commit1} \textit{commit2} & zeigt Unterschiede zwischen zwei Commits
  \end{tabu}
\end{frame}

\begin{frame}{\texttt{git commit}}
  \begin{tabu}{>{\ttfamily}l X[,L]}
    git commit                       & erzeugt Commit aus jetzigem Staging-Bereich, öffnet Editor für Commit-Message \\
    git commit -m "\textit{message}" & Commit mit \texttt{\textit{message}} als Message \\
    git commit --amend               & letzten Commit ändern (fügt aktuellen Staging hinzu, Message bearbeitbar) \\
    & \alert{\bfseries Niemals commits ändern, die schon gepusht sind!}
  \end{tabu}

  \begin{itemize}
    \item Wichtig: Sinnvolle Commit-Messages
      \begin{itemize}
        \item Erster Satz ist Zusammenfassung (ideal < 50 Zeichen)
        \item Danach eine leere Zeile lassen
        \item Dann längere Erläuterung des commits
      \end{itemize}
    \item Logische Commits erstellen, für jede logische Einheit ein Commit
      \begin{itemize}
        \item \texttt{git add -p} ist hier nützlich
      \end{itemize}
    \item Hochgeladene Commits sollte man nicht mehr ändern
  \end{itemize}
\end{frame}

\begin{frame}{\texttt{git pull}, \texttt{git push}}
  \begin{tabu}{>{\ttfamily}l X[,L]}
    git pull          & Commits herunterladen \\
    git push          & Commits hochladen
  \end{tabu}

  \begin{itemize}
    \item Aus der Installationsanleitung: \\
      \texttt{git config --global pull.rebase true}
  \end{itemize}
\end{frame}

\begin{frame}[fragile]{Achtung: Merge conflicts}
  \begin{center}
    \huge Don't Panic
  \end{center}

  Entstehen, wenn \texttt{git} nicht automatisch mergen kann (selbe Zeile geändert, etc.)

  \begin{enumerate}
    \item Die betroffenen Dateien öffnen
    \item Markierungen finden und die Stelle selbst mergen (meist wenige Zeilen)
      \begin{verbatim}
        <<<<<<< HEAD
        foo
        ||||||| merged common ancestors
        bar
        =======
        baz
        >>>>>>> Commit-Message
\end{verbatim}
    \item Merge abschließen:
      \begin{enumerate}
        \item \texttt{git add …}
        \item \texttt{git rebase --continue}
      \end{enumerate}
  \end{enumerate}
  Nützlich: \texttt{git config --global merge.conflictstyle diff3}
\end{frame}

\begin{frame}{\texttt{git checkout}}
  \begin{tabu}{>{\ttfamily}l X[,L]}
    git checkout \textit{commit} & Commit ins Arbeitsverzeichnis laden \\
    git checkout \textit{file}   & Änderungen an Dateien verwerfen (zum letzten Commit zurückkehren)
  \end{tabu}
\end{frame}

\begin{frame}{\texttt{git stash}}
  \begin{tabu}{>{\ttfamily}l X[,L]}
    git stash     & Änderungen kurz zur Seite schieben \\
    git stash pop & Änderungen zurückholen aus Stash
  \end{tabu}
\end{frame}

\begin{frame}[fragile]{\texttt{.gitignore}}
    \begin{itemize}
    \item Man möchte nicht alle Dateien von \texttt{git} beobachten lassen
    \item z.B. \texttt{build}-Ordner
    \end{itemize}
    \begin{center}
        \Large Lösung: \texttt{.gitignore}-Datei
    \end{center}

    \begin{itemize}
    \item einfache Textdatei
    \item enthält Regeln für Dateien, die nicht beobachtet werden sollen
    \end{itemize}
    Beispiel:
    \vspace{1em}
    \begin{lstlisting}[language=, keywordstyle={}]
      build/
      *.pdf
      __pycache__/
    \end{lstlisting}
\end{frame}

\begin{frame}{Hoster}
  \begin{tabu}{@{} X[,C] @{} X[,C] @{} X[,C] @{}}
    \href{https://github.com}{\includegraphics[width=0.75\linewidth]{figures/github.png}} &
    \href{https://bitbucket.org}{\includegraphics[width=0.75\linewidth]{figures/bitbucket.png}} &
    \href{https://gitlab.com}{\includegraphics[width=0.75\linewidth]{figures/gitlab.png}} \\
    \begin{itemize}
      \item größter Hoster
      \item viele open-source Projekte
      \item Unbegrenzt private Repositories für Studenten und Forscher:  \newline
        \href{http://education.github.com}{education.github.com}
    \end{itemize}
    &
    \begin{itemize}
      \item kostenlose private Repos mit höchstens fünf Leuten
      \item keine Speicherbegrenzungen
      \item Hängt was Oberfläche und Funktionen angeht, den beiden anderen weit hinterher
    \end{itemize}
    &
    \begin{itemize}
      \item open-source
      \item keine Begrenzungen an privaten Repos
      \item kann man selbst auf einem eigenen Server betreiben
    \end{itemize}
  \end{tabu}
  \begin{center}
    \onslide<2->{%
      \Large%
      \enquote{Now, everybody sort of gets born with a GitHub account}
      – Guido van Rossum
    }
  \end{center}
\end{frame}

\begin{frame}{SSH-Keys}
  Git kann auf mehrere Arten mit einem Server kommunizieren:
  \begin{itemize}
    \item HTTPS: funktioniert immer, keine Einstellungen erforderlich, Passwort muss für jede Kommunikation eingegeben werden
    \item SSH: Keys müssen erzeugt und eingestellt werden, Passwort für den Key muss nur einmal pro Session eingegeben werden.
  \end{itemize}

  SSH-Keys:
  \begin{enumerate}
    \item \texttt{ssh-keygen}
    \item Passwort wählen
    \item \texttt{cat \textasciitilde/.ssh/id\_rsa.pub}
    \item Ausgabe ist Public-Key, beim Server eintragen (im Browser)
  \end{enumerate}
\end{frame}
